\chapter*{Conclusiones}
\addcontentsline{toc}{chapter}{Conclusiones}
\markboth{\MakeUppercase{Conclusiones}}{\MakeUppercase{Conclusiones}}

El objetivo de este trabajo fue desarrollar una biblioteca en Python que permitiera ejecutar operaciones tensoriales e inferencia de redes neuronales sobre datos cifrados con el esquema de cifrado homomórfico CKKS, aprovechando la capacidad de procesamiento paralelo de la GPU. La biblioteca desarrollada cumple ese objetivo: permite describir una red neuronal de forma declarativa o convertir un modelo previamente entrenado en PyTorch, compilarla contra un contexto criptográfico y ejecutar su inferencia sobre datos cifrados, sin que el usuario manipule de forma directa las claves, los niveles de la cadena de módulos ni los objetos del \textit{backend}.

Del diseño y la implementación descritos en el Capítulo 4 se desprenden varias conclusiones. La primera es que la abstracción del esquema criptográfico es viable sin sacrificar funcionalidad: la totalidad de la complejidad de CKKS, incluida la generación de claves, la relinealización, el reescalado y el \textit{bootstrapping}, pudo ocultarse detrás de una interfaz equivalente a la de las bibliotecas de aprendizaje profundo de uso común. La excepción son los parámetros del esquema, que permanecen expuestos al usuario porque su selección automática en función de la arquitectura de la red no resultó abordable dentro del alcance de este trabajo.

La segunda conclusión se refiere a la gestión del presupuesto de ruido. Determinar durante la compilación la profundidad multiplicativa total del modelo y contrastarla con los niveles disponibles en el contexto permite decidir automáticamente si la red requiere \textit{bootstrapping}, y evitar su costo cuando el circuito cabe en la cadena de módulos. Complementar esa decisión con un refresco disparado automáticamente en tiempo de inferencia, incluso en medio de la evaluación de una activación, libera al usuario de una de las tareas más delicadas del cómputo homomórfico.

La tercera conclusión es que la calibración de los rangos de las activaciones no es una optimización opcional sino un requisito de corrección. La aproximación polinómica de ReLU solo es precisa mientras su entrada permanece dentro de un intervalo acotado, y fuera de él el error crece con rapidez suficiente para invalidar el resultado de la inferencia. Plegar los rangos estimados en los pesos en texto plano durante la compilación resuelve este problema sin introducir ninguna operación homomórfica adicional en tiempo de inferencia.

La cuarta conclusión es que el conjunto de pruebas automatizadas construido junto con la biblioteca resultó ser el principal instrumento de verificación del trabajo, al constituir la garantía de que cada operación cifrada reproduce, dentro del margen de error propio de la aritmética aproximada de CKKS, el resultado del cómputo equivalente en texto plano.

De la evaluación comparativa se desprenden otras tres conclusiones, de carácter cuantitativo. La primera es que la aceleración por GPU cumple su propósito. La biblioteca resultó la implementación cifrada más rápida en los dos casos de uso: resuelve una inferencia de la red de regresión en 6.68 segundos, frente a los 78.24 de Orion y los 790.03 de Concrete~ML, y una de la red convolucional en 28.31 segundos, frente a 46.49 y 448.88. Esto supone aceleraciones de 11.7 y 1.6 veces sobre Orion, que emplea el mismo esquema sobre CPU, y de 118 y 15.9 veces sobre Concrete~ML. La ventaja se estrecha en la red convolucional porque su costo se concentra en rotaciones, que se benefician menos de la paralelización que las multiplicaciones y los reescalados de las activaciones.

La segunda es que esa ventaja se paga en memoria, y que el conjunto de claves de Galois es el factor dominante. En la red de regresión, con las claves alojadas en la memoria de video, el conjunto de trabajo alcanzó 17.6~GB; en la red convolucional, cuyo empaquetado exige un conjunto de claves mucho mayor, hubo que trasladarlas a la memoria del anfitrión, lo que redujo el consumo de video a 3.0~GB a cambio de 43.6~GB de memoria principal. Este intercambio es lo que permitió evaluar la red convolucional en una tarjeta de 24~GB, y es también el factor que acota el tamaño de las redes que la biblioteca puede procesar.

La tercera es que la fidelidad de la inferencia cifrada depende de la arquitectura evaluada, y que su determinante es el número de operaciones homomórficas que el circuito acumula. En la red de regresión la biblioteca reprodujo la salida en texto plano con un error absoluto medio de $6.2 \times 10^{-4}$, el menor de las tres implementaciones evaluadas. En la red convolucional, en cambio, el error sobre las puntuaciones ascendió a 91.2 y la clase predicha difirió de la del modelo en claro en el 40\% de las muestras, mientras las dos bibliotecas de referencia preservaron íntegramente la decisión del clasificador. La causa es la implementación de la capa convolucional: expresarla como una multiplicación matriz-vector de propósito general otorga flexibilidad sobre los parámetros de la capa y reutiliza una única ruta de cómputo, pero a costa de un número de operaciones mayor que el de un empaquetado específico, y en CKKS cada operación añade ruido al texto cifrado. Esa acumulación no se manifiesta al ejecutar el mismo modelo aproximado sin cifrar, donde la aritmética de punto flotante no arrastra el error, lo que confirma que el problema pertenece al cómputo cifrado y no a la aproximación de las activaciones.

\section*{Limitaciones}
\addcontentsline{toc}{section}{Limitaciones}

El trabajo presenta las siguientes limitaciones. La más relevante es la señalada por la evaluación: la implementación actual de la convolución no preserva la fidelidad de la inferencia en la red convolucional evaluada, de modo que la biblioteca resulta utilizable para redes totalmente conectadas pero no, en su estado actual, para arquitecturas convolucionales.

La selección de los parámetros criptográficos permanece a cargo del usuario, quien debe conocer la profundidad multiplicativa de su red para elegir una cadena de módulos adecuada. El tamaño de los modelos que pueden evaluarse está acotado por dos factores: la profundidad multiplicativa del circuito homomórfico y el consumo de memoria de video, del que el conjunto de claves de Galois representa la mayor parte cuando se emplea un grado de polinomio elevado. A ello se añade que el gestor de memoria del \textit{backend} reserva por adelantado prácticamente la totalidad de la tarjeta, con independencia del tamaño de la red, lo que impide compartir el dispositivo con otro proceso.

El catálogo de capas soportadas se limita a las capas lineal y convolucional junto con la activación ReLU, sin incluir agrupamiento, normalización por lotes ni conexiones residuales. La biblioteca cubre únicamente la inferencia y no el entrenamiento sobre datos cifrados. Por último, la calidad de la aproximación de las activaciones depende de que el conjunto de datos empleado en la calibración sea representativo del rango de operación real de la red.

\section*{Recomendaciones}
\addcontentsline{toc}{section}{Recomendaciones}

A partir de las limitaciones anteriores se identifican varias líneas de continuación, ordenadas según la prioridad que sugiere la evaluación.

La primera es rediseñar la implementación de la capa convolucional. La evaluación mostró que el número de operaciones homomórficas que acumula es el origen de la pérdida de fidelidad, y también la razón de que la ventaja en tiempo frente a las implementaciones sobre CPU se estreche en esta arquitectura. Adoptar un empaquetado específico de la convolución, que reordene el cómputo para obtener el mismo resultado con menos operaciones, corregiría ambos problemas. El compromiso a evaluar es la pérdida de flexibilidad sobre los parámetros de la capa, pues los esquemas más eficientes admiten un conjunto más restringido de configuraciones; una solución posible es conservar la ruta general como alternativa para los casos que el empaquetado especializado no cubra.

La segunda es buscar una forma más eficiente de evaluar las funciones de activación. El perfilado por capa las señala como la operación más costosa del circuito, con el 94\% del tiempo en la red de regresión, y su profundidad multiplicativa es la que obliga a recurrir al \textit{bootstrapping}. Las tablas de búsqueda constituyen una alternativa a la aproximación polinómica, ya que evalúan la función de forma exacta y con una profundidad independiente de la precisión buscada. Dado que CKKS no las soporta de forma nativa, explorar esta vía supone considerar esquemas que sí lo hagan, como TFHE, ya sea sustituyendo el esquema o combinando ambos dentro de la misma inferencia mediante conversión entre representaciones.

La tercera es intervenir sobre el \textit{backend}. Varias de las limitaciones observadas no pertenecen a la biblioteca sino a HEonGPU: la reserva anticipada de la totalidad de la memoria de video, la ausencia de tablas de búsqueda y el costo del conjunto de claves de Galois. Mantener una bifurcación propia del proyecto, o contribuir directamente a su desarrollo, permitiría incorporar estas capacidades y una gestión de memoria más granular, y beneficiaría además a la comunidad que emplea esa biblioteca.

A estas se suman tres líneas de menor urgencia. La selección automática de los parámetros del esquema a partir de la profundidad multiplicativa y de los rangos estimados de la red eliminaría el último punto en que el usuario debe razonar en términos criptográficos. La ampliación del catálogo de capas admitiría arquitecturas más profundas y de uso más extendido. Y el empaquetado de los datos en múltiples textos cifrados permitiría evaluar redes cuyas capas exceden el número de ranuras disponibles en un único texto cifrado.

\section*{Aporte}
\addcontentsline{toc}{section}{Aporte}

El aporte principal de este trabajo es la biblioteca resultante, publicada como paquete instalable de código abierto y disponible para la comunidad. Al ofrecer inferencia privada acelerada por GPU detrás de una interfaz familiar para quienes trabajan con aprendizaje profundo, la biblioteca reduce la barrera de entrada al cómputo sobre datos cifrados para usuarios sin formación especializada en criptografía. A ello se añade el banco de pruebas construido para la evaluación, que ejecuta cada implementación en un proceso aislado y mide de forma homogénea el tiempo, la memoria y la fidelidad de bibliotecas basadas en esquemas distintos, y que constituye por sí mismo un instrumento reutilizable para comparar futuras implementaciones de inferencia sobre datos cifrados.
